% alpha_intro.tex — Abstract and Introduction
% "The Economics of Extended Play: What Happens When a Regulator Adds Time to the Clock?"
% Draft: alpha

\begin{abstract}
When a regulator extends the duration of a high-stakes competitive interaction, does total output increase, or do participants merely redistribute effort across the expanded time horizon? We study this question using FIFA's 2022 directive instructing referees to increase stoppage time in professional football, a voluntary enforcement shift that generated differential compliance across leagues. Exploiting variation in enforcement intensity across 15 leagues over 8 seasons (36,267 matches), our preferred difference-in-differences specification estimates that treated leagues increased second-half stoppage time by 0.744 minutes relative to controls (wild cluster bootstrap $p = 0.036$). Within-country league pairs sharing the same refereeing infrastructure yield a supplementary estimate of 0.868 minutes (WCB $p = 0.074$). Five findings emerge. First, compliance was near-universal: 100\% of referees with 20 or more post-directive matches increased mean stoppage time, with 87\% match-level compliance in the first five matchweeks. Second, additional time reshuffled goals rather than creating them --- goals scored after the 90th minute rose by 0.030 per match ($p < 0.001$), but total goals per match were unchanged ($\beta = -0.033$, $p = 0.39$), as the per-minute scoring rate fell from 0.72 to 0.50 ($t = 19.9$, $p < 0.001$). Third, betting market accuracy was unaffected (Pinnacle Brier score DiD: $+0.007$, $p = 0.41$). Fourth, game-changing goals scored after the 85th minute increased by 0.035 per match ($p < 0.001$). Fifth, player injuries did not increase despite 0.76 additional minutes of workload per player per match; total injuries fell modestly ($-0.09$, $p = 0.005$) and muscular injuries were unchanged ($p = 0.73$, 88.6\% power for a 10\% increase). We are transparent about identification: the event study does not support parallel pre-trends ($F = 9.0$, $p < 0.001$) and the \citet{RambachanRoth2023} honest DiD breakdown value is $\bar{M}^* = 0$, so our causal interpretation rests on institutional arguments rather than statistical pre-trend tests.
\end{abstract}

\section{Introduction}\label{sec:introduction}

Time is a scarce resource in competitive environments. When a regulator expands the time available for a contest---whether by lengthening trading hours in a financial market, extending a legislative session, or adding minutes to an athletic competition---the economic consequences depend on whether participants treat the additional time as an opportunity for new productive activity or as a cushion that dilutes effort per unit of time. The distinction matters for policy: if extending duration increases total output, regulators can raise the stakes of competition at low cost; if it merely redistributes existing output across a longer window, the intervention is a transfer from the intensity margin to the quantity margin with ambiguous welfare implications. This tradeoff connects to foundational questions in labor economics about the relationship between hours and productivity \citep{Pencavel2015} and to market design research on how closing rules shape bidding behavior \citep{RothOckenfels2002}. Yet credible causal evidence on time expansion is scarce, because regulators rarely change the duration of a competitive interaction in a discrete, well-documented, and plausibly exogenous fashion.

This paper exploits a rare natural experiment: FIFA's 2022 directive to referees instructing them to enforce additional time more generously in professional football (soccer). In September 2022, Pierluigi Collina, chairman of FIFA's Referees Committee, issued guidance ahead of the Qatar World Cup directing referees to add time more accurately for stoppages during play. The directive was not a rule change---the Laws of the Game, maintained by the International Football Association Board (IFAB), already granted referees discretion over stoppage time. Rather, it was an enforcement shift: a voluntary recalibration of existing discretion, amplified by the high-profile World Cup matches in which stoppage time routinely exceeded 10 minutes. IFAB endorsed the approach at its March 2023 Annual General Meeting using advisory language (``should'' rather than ``must''), and UEFA explicitly declined to mandate the directive for the Champions League and Europa League. The result was a natural enforcement gradient: top domestic leagues in England, Spain, Germany, Italy, and France adopted the directive with varying speed and intensity, while lower divisions and leagues in other countries responded more modestly.

Our identification strategy exploits this enforcement gradient rather than a binary treatment. No league was fully untreated---all 15 leagues in our sample increased stoppage time to some degree, with control leagues averaging a 0.89-minute increase compared to 1.57 minutes in the Big~5 leagues. The identifying variation comes from the \emph{differential} intensity of enforcement across leagues that were otherwise subject to the same underlying rule. We pursue two complementary strategies. First, our preferred specification pools all 15 leagues (36,267 matches across 8 seasons) and estimates the differential increase via a standard two-period difference-in-differences design, clustering at the league level and conducting inference with the wild cluster bootstrap appropriate for 15 clusters \citep{CameronGelbachMiller2008}. Second, we construct within-country league pairs---English Premier League and Championship, German Bundesliga and 2.~Bundesliga, Italian Serie~A and Serie~B, French Ligue~1 and Ligue~2---that share the same national refereeing infrastructure, providing a tighter comparison that holds referee labor markets constant. The within-country estimate ($\beta = 0.868$, WCB $p = 0.074$) is larger in magnitude than the all-controls estimate ($\beta = 0.744$, WCB $p = 0.036$), consistent with the interpretation that controlling for shared referee pools sharpens the contrast.

Our paper makes five contributions:
\begin{enumerate}
    \item \textbf{Time expansion and the intensive margin.} We provide, to our knowledge, the first causal evidence on the quantity--intensity decomposition of a time expansion in a competitive setting. A Gelbach decomposition reveals that 87\% of the increase in late goals operates through the quantity channel (more minutes of play), while the per-minute scoring rate is unchanged ($t = 0.41$, $p = 0.68$ for goals after the 90th minute). For total goals, the per-minute rate \emph{fell} from 0.72 to 0.50 ($t = 19.9$, $p < 0.001$), exactly offsetting the additional minutes.
    \item \textbf{Compliance and institutional design.} We document near-universal compliance with a voluntary directive: 100\% of the 85 referees with 20 or more post-directive matches increased their mean stoppage time, and match-level compliance reached 70.2\% overall and 87\% within the first five matchweeks. This speaks to the power of focal-point coordination among agents who share professional norms \citep{Schelling1960}.
    \item \textbf{Betting market efficiency.} Despite a large shift in match structure, bookmaker forecast accuracy was unchanged, with Brier score differences of $+0.007$ ($p = 0.41$) for Pinnacle and $+0.005$ ($p = 0.58$) for Bet365. Markets absorbed the regime change without measurable loss of calibration.
    \item \textbf{Worker health.} Additional playing time raised per-player workload by 0.76 minutes per match ($p < 0.001$), yet total injuries decreased ($-0.09$, $p = 0.005$) and muscular injuries---the category most plausibly linked to fatigue---were unchanged ($-0.007$, $p = 0.73$), with our test well-powered to detect a 10\% increase (88.6\% power).
    \item \textbf{Game-changing moments.} Goals scored after the 85th minute in matches where the score was close increased by 0.035 per match ($p < 0.001$), suggesting that extended time raises the drama and competitive tension of matches without increasing total scoring.
\end{enumerate}

The magnitudes are robust across specifications. A specification curve spanning 420 alternative models yields a median coefficient of 0.669, with 88.8\% of estimates significant at the 5\% level and 95\% carrying a positive sign. Leave-one-out analysis dropping each league in turn produces a range of $[0.587, 0.972]$, all individually significant. A permutation test shuffling treatment assignment across leagues returns a two-sided $p$-value of 0.045. The Big~5-only estimate, which compares the five most intensely treated leagues before and after the directive, is 1.575 minutes ($\text{SE} = 0.250$, $p < 0.001$). Excluding the 2022--23 transitional season strengthens the estimate to $\beta = 0.810$.

We are equally forthcoming about the limitations of our design. The event study rejects parallel pre-trends ($F = 9.0$, $p < 0.001$), and the \citet{RambachanRoth2023} honest difference-in-differences sensitivity analysis yields a breakdown value of $\bar{M}^* = 0$, meaning the result is not robust to any extrapolation of pre-existing differential trends. This is the single most important caveat for causal interpretation. We argue that the institutional context---a well-documented, discrete directive with observable compliance that was differentially enforced for reasons unrelated to league-specific trends in stoppage time---provides a credible basis for a causal reading, but we do not ask the reader to accept this on the strength of the pre-trends alone. Supporting this interpretation, a first-half goals placebo test shows no differential change in first-half scoring ($p = 0.33$), confirming that the treatment affected second-half stoppage time specifically rather than reflecting a general divergence in match characteristics.

This paper contributes to several literatures. Most directly, it extends the economics of sports regulation, where \citet{GaricanoEtAl2005} showed that referees strategically manipulate stoppage time to favor home teams, and \citet{ButlerEtAl2026} studied strategic behavior in added time. Our contribution differs in studying the consequences of a \emph{regime change} in stoppage time norms rather than strategic manipulation within an existing regime, and in combining difference-in-differences identification across 15 leagues with betting market, compliance, and health outcome data. \citet{Kocsoy2025} provides a concurrent analysis, but does not exploit within-country variation, conduct wild cluster bootstrap inference, or examine the quantity--intensity decomposition. More broadly, our findings inform the market design literature on time-extension rules \citep{RothOckenfels2002} by providing field evidence that extending the clock shifts the timing of decisive events without expanding total output---a pattern consistent with the ``sniping'' equilibrium in hard-close auctions, where last-minute activity intensifies at the boundary rather than spreading across the expanded window. Finally, our results contribute to the labor economics literature on hours and productivity \citep{Pencavel2015,GoldmanSeskin1983} by documenting a setting in which extending the work period reduced output per unit of time while leaving total output unchanged.

The remainder of the paper proceeds as follows. Section~\ref{sec:background} describes the institutional setting, the FIFA directive, and the enforcement gradient across leagues. Section~\ref{sec:data} introduces our data, covering match-level stoppage times, goals, betting odds, referee assignments, and injury records across 15 leagues and 8 seasons. Section~\ref{sec:empirical} presents the empirical strategy, including the difference-in-differences design, wild cluster bootstrap inference, and the specification curve and permutation tests. Section~\ref{sec:results} reports results on stoppage time, goals, betting markets, compliance, and worker health. Section~\ref{sec:robustness} conducts robustness analysis, including the event study, honest DiD sensitivity, leave-one-out tests, and placebo checks. Section~\ref{sec:conclusion} concludes.
